% Slide — Novas: Zabraoui
\begin{frame}{\novo: três heurísticas do artigo-base (Zabraoui et al., 2025)}
  \footnotesize Seção 3.5 do artigo-base define quatro políticas de
  referência; apenas $(s,S)$ já existia no portfólio. As outras três entram
  aqui.

  \renewcommand{\arraystretch}{0.95}\scriptsize
  \begin{tabular}{@{}p{0.24\linewidth}p{0.66\linewidth}@{}}
    \toprule
    \textbf{Política} & \textbf{Regra adotada} \\
    \midrule
    Min-Max &
      Pede até $M$ quando posição $\leq m$;
      $m=\mu L$ (lead time $L$), $M=m + \mu\cdot$cobertura \\
    Fixed Interval Repl. &
      Revisão a cada $R$ ciclos; repõe até
      $S=\mu(R{+}L)+z\sigma\sqrt{R{+}L}$ \\
    Vendor-Responsive &
      $s_t=\mu L + z_t\sigma\sqrt{L}$,
      $z_t=z_0(1+\kappa(\alpha_{\text{alvo}}-\mathrm{SL}_t))$ -- realimentação
      pelo serviço realizado \\
    \bottomrule
  \end{tabular}
  \vspace{0.3em}

  \begin{alertblock}{Ressalva de fidelidade, registrada no código}
    \scriptsize Min-Max e Fixed Interval: descritas no artigo em
    \textbf{uma frase cada, sem fórmula}. Vendor-Responsive é a mais
    subespecificada (``\textit{adapts order timing based on supplier lead
    times and historical service levels}''). Fórmulas acima são leitura
    razoável, \highlight{não transcrição literal} -- dito explicitamente
    nos comentários do código.
  \end{alertblock}

  \note{
    Este slide existe para ser transparente sobre o grau de fidelidade de
    cada implementação. O artigo do Zabraoui é a referência que este
    trabalho adota para o que não tem estado da arte próprio -- mas o
    artigo descreve essas três políticas em uma frase cada, sem equação.
    Documentei no próprio código, com essa mesma ressalva, qual foi a
    leitura escolhida: Min-Max e Fixed Interval seguem a definição
    clássica de manual de operações; Vendor-Responsive foi a mais livre de
    interpretar, porque "adapta o timing do pedido com base no lead time
    do fornecedor e no nível de serviço histórico" não tem parâmetro
    algum -- montei um laço de realimentação em que o z-score de segurança
    sobe quando o serviço realizado fica abaixo da meta. É importante que
    a banca saiba que isso é uma escolha nossa razoável, não uma
    reprodução certificada do artigo.
  }
\end{frame}
